\documentclass[12pt]{article}
\usepackage[T1]{fontenc}
\usepackage[utf8]{inputenc}
\usepackage{lmodern}
\usepackage{textcomp}
\usepackage{enumitem}
\usepackage{geometry}
\geometry{margin=1in}
\begin{document}
\begin{center}
Answer Key - Geography - K-12 Level Assessment 1 \
\small (Answers are ordered exactly as in the question paper.)
\end{center}

\section*{Multiple Choice}
\begin{enumerate}[label=\arabic*.]
\item \textbf{Answer:} D
\\ \textit{Options:} A) Longitude; B) Latitude; C) Equator; D) Distance to the nearest city
\\ \textit{Note:} Distance to the nearest city is not used to determine absolute location because absolute location is defined by specific coordinates (latitude and longitude) rather than relative measures like distance to a reference point.
\item \textbf{Answer:} C
\\ \textit{Options:} A) A central park; B) Prominent religious buildings; C) A large central square surrounded by government and administrative buildings; D) Luxury apartment buildings
\\ \textit{Note:} A large central square surrounded by government and administrative buildings is a distinct characteristic of East European cities, reflecting their historical importance as centers of political power and public life.
\item \textbf{Answer:} B
\\ \textit{Options:} A) capitals of countries formerly colonized by the English; B) destinations for vast numbers of pilgrims; C) financial centers for a large fraction of the world's economy; D) examples of modern urban planning
\\ \textit{Note:} Both Varanasi and Mecca serve as significant religious sites that attract millions of pilgrims annually, making them vital destinations for spiritual journeys in their respective faiths.
\item \textbf{Answer:} D
\\ \textit{Options:} A) new residents tear down and rebuild housing units.; B) new residents live peacefully with old residents.; C) new residents assimilate with old residents completely.; D) new residents move into areas occupied by older resident groups.
\\ \textit{Note:} The correct answer highlights the key concept of ecological succession, where new species or groups migrate into an area, leading to changes in the existing community dynamics established by older resident groups.
\item \textbf{Answer:} B
\\ \textit{Options:} A) regionalist groups.; B) separatist groups.; C) terrorist groups.; D) interventionist groups.
\\ \textit{Note:} The Basques, Bretons, Kashmiris, and Tamils are considered separatist groups because they each seek greater autonomy or independence for their distinct cultural, ethnic, or linguistic identities within larger nation-states.
\end{enumerate}

\section*{True / False}
\begin{enumerate}[label=\arabic*.]
\setcounter{enumi}{5}
\item \textbf{Answer:} True
\\ \textit{Options:} A) True; B) False
\\ \textit{Note:} Hinduism is considered one of the oldest religions because its roots trace back to ancient practices and beliefs in the Indian subcontinent, with some of its texts and traditions dating back over 4,000 years.
\item \textbf{Answer:} True
\\ \textit{Options:} A) True; B) False
\\ \textit{Note:} Students often use mental maps to efficiently recall and navigate familiar routes, drawing on their spatial knowledge and experiences to reach school each day.
\item \textbf{Answer:} True
\\ \textit{Options:} A) True; B) False
\\ \textit{Note:} China's establishment of special economic zones (SEZs) has successfully attracted foreign investment and stimulated economic growth by offering favorable conditions such as tax incentives and less regulatory oversight.
\item \textbf{Answer:} True
\\ \textit{Options:} A) True; B) False
\\ \textit{Note:} Guyana, Suriname, and French Guiana are all situated on the northeastern coast of South America, making them part of the South American continent.
\item \textbf{Answer:} True
\\ \textit{Options:} A) True; B) False
\\ \textit{Note:} The statement is true because a Mercedes-Benz dealership, as a luxury car retailer, caters to a more affluent customer base and serves a broader geographic region compared to standard car dealerships, reflecting its higher-order central place function in the market hierarchy.
\end{enumerate}

\section*{Long Form Response}
\begin{enumerate}[label=\arabic*.]
\setcounter{enumi}{10}
\item \textbf{Answer:} The rank-size rule is a principle in urban geography that describes how the population of cities within a region is distributed in relation to their rank in the urban hierarchy. According to this model, the population of a city is inversely proportional to its rank; for example, the second-largest city will have about half the population of the largest city, the third-largest will have one-third, and so on. This pattern suggests that as cities decrease in size, their population declines in a predictable manner. For instance, in the United States, New York City is the largest urban center, followed by Los Angeles, and then Chicago, illustrating this relationship. The rank-size rule helps explain urban organization and can reveal economic and social dynamics within a region.
\\ \textit{Note:} The answer correctly describes the rank-size rule by illustrating how the population of cities decreases in a predictable manner based on their rank within the urban hierarchy, providing a clear understanding of the inverse relationship between city rank and population size.
\item \textbf{Answer:} Export processing zones (EPZs) are designated areas within a country where goods can be manufactured and exported with minimal customs regulations and taxes. One key feature of EPZs is that they often attract foreign investment by providing incentives such as tax breaks and streamlined procedures. However, the environmental regulations in these zones can be less strict compared to those in other areas, which can lead to challenges in pollution management. This leniency may encourage industries to operate with more freedom, potentially resulting in higher levels of pollution if proper measures are not enforced. Therefore, while EPZs can boost economic growth, their relaxed environmental regulations can have negative implications for local ecosystems and public health if not managed carefully.
\\ \textit{Note:} correct because it accurately describes the defining characteristics of export processing zones, particularly their focus on attracting foreign investment through incentives, while also highlighting the potential for weaker environmental regulations that can complicate pollution management efforts.
\end{enumerate}

\vfill
\noindent\textit{End of Answer Key}
\end{document}